\section{Recognize complete byte patterns}
A continuation byte has high bits \texttt{10}, giving an unsigned range of
hexadecimal \texttt{80} through \texttt{bf}. Leading-byte classes determine the
required number of continuations. Four leading-byte values impose an additional
restriction on the first continuation. The ranges below follow Unicode 17's
UTF-8 definition.\footnote{\href{https://www.unicode.org/versions/Unicode17.0.0/core-spec/chapter-3/}{Unicode 17, Chapter 3, Section 3.9.3, Table 3-7}. See also \href{https://www.rfc-editor.org/rfc/rfc3629}{RFC 3629, Section 4}, for a byte grammar.}
All byte values in the table are hexadecimal. A dash means that no byte occupies
that position.

\Needspace{16\baselineskip}
\begin{longtable}{llll}
  \toprule Lead   & Next byte       & Third byte      & Fourth byte     \\\midrule
  \texttt{00--7f} & ---             & ---             & ---             \\
  \texttt{c2--df} & \texttt{80--bf} & ---             & ---             \\
  \texttt{e0}     & \texttt{a0--bf} & \texttt{80--bf} & ---             \\
  \texttt{e1--ec} & \texttt{80--bf} & \texttt{80--bf} & ---             \\
  \texttt{ed}     & \texttt{80--9f} & \texttt{80--bf} & ---             \\
  \texttt{ee--ef} & \texttt{80--bf} & \texttt{80--bf} & ---             \\
  \texttt{f0}     & \texttt{90--bf} & \texttt{80--bf} & \texttt{80--bf} \\
  \texttt{f1--f3} & \texttt{80--bf} & \texttt{80--bf} & \texttt{80--bf} \\
  \texttt{f4}     & \texttt{80--8f} & \texttt{80--bf} & \texttt{80--bf} \\\bottomrule
\end{longtable}
A leading-byte classifier must reject a standalone continuation, \texttt{c0} or
\texttt{c1}, and \texttt{f5} through \texttt{ff}. It must then enforce both bounds
on each continuation. Classifying only the number of high leading bits and accepting
any \texttt{80--bf} tail would admit encodings outside this table.

\begin{MARSWorkedExample}{The value behind three bytes}
  For \texttt{e2 82 ac}, remove the structural leading bits. The remaining values
  are hexadecimal 2, 2, and 2c, or decimal 2, 2, and 44. Their bit positions give
  \[
    2\cdot2^{12}+2\cdot2^6+44
    =8192+128+44=8364=\texttt{20ac}_{16}.
  \]
  The scanner can validate and count this encoding without constructing 8364 in a
  register. Retaining the expected continuation count and next-byte bounds is enough
  for its interface. A decoder returning scalar values would need additional output
  and accumulation behavior.
\end{MARSWorkedExample}
